\begin{table*}[!tbp]
\centering\scriptsize\setlength\tabcolsep{3pt}\renewcommand{\arraystretch}{1.0}
\caption{A minimum audit workflow for one attribute $c$ (a finding, a demographic attribute or an acquisition variable) in a frozen encoder or a medical multimodal LLM, assembled from the controls that \S\ref{sec:evidence} and \S\ref{sec:evaluation} of the article identify as necessary. Each step names its inputs, the unit at which data are split, the null or control that must accompany the measurement, the output to report and the criterion for proceeding; the tooling column points to the software of \S\ref{sec:eval-bench}. Steps 1--2 are feasible with embedding-only access, steps 3--4 need activation access to open weights; a null at any step is a reportable outcome, and the workflow measures detectability, not clinical usefulness or a deployment decision. The workflow is a recommendation of the authors, not a coded finding.}
\label{tab:workflow}
\begin{tabular}{@{}L{1.6cm}L{2.6cm}L{1.45cm}L{2.95cm}L{2.6cm}L{1.95cm}L{1.8cm}@{}}
\toprule
Step & Inputs & Split unit & Null or control & Output reported & Stronger evidence when & Tooling \\
\midrule
1. Availability of $c$ & frozen embeddings at the stated locus and pooling (Table~\ref{tab:models}); labels for $c$; a comparison encoder & patient, slide or subject; an unused cohort held out & control task or permuted labels; randomly initialised encoder of the same architecture; general encoder on the same data & AUROC or AUC with intervals over resamples and seeds; task-minus-control gap; layer or pooling at which $c$ first appears & $c$ decodable above both nulls on the held-out cohort & scikit-learn or the probing code of the cited studies; DINOv2/SigLIP checkpoints for the general comparator \\
2. Where along the path & the same probe class, capacity and tuning at encoder output, connector output and successive language-model layers & as step 1 & the same nulls at every locus & accuracy profile over loci with intervals; the locus of the largest drop & a drop that exceeds interval width at a named locus & TransformerLens or nnsight hooks on the open MLLMs of Table~\ref{tab:models} \\
3. Direction and controls & $v_c$ from the step-1 probe (or a validated dictionary feature); $K$ matched control directions ($K$ fixed in advance; 99 and 199 are illustrative values chosen for Monte Carlo resolution, not validated requirements: norm- and spectrum-matched; unrelated-concept; label-permuted, refit through the same pipeline) & held-out inputs disjoint from the data used to fit $v_c$ & the direction-selectivity screen of \S\ref{sec:future}; both signs; pre-specified $\alpha$ grid & effect of $v_c$ against each control distribution, per magnitude and sign, with intervals over inputs and constructed directions & effect exceeds the 95th percentile of the label-permuted family (the inferential control) for a pre-registered primary cell or as a maximum over the grid, with patient-level resampling; the other two quantiles reported as descriptive checks & pyvene or nnsight steering hooks; LEACE for erasure \\
4. Behavioural end-point & the model's own output (report, answer, class score); an off-target end-point & as step 3 & unrelated-concept direction; sham intervention & change in the pre-specified end-point and in the off-target end-point & on-target change with no off-target change beyond noise & the model's generation code; RadFact or CheXbert for report end-points \\
5. Reporting & everything above & -- & -- & the reporting card of \S\ref{sec:future} (locus, layer, tokens, construction data, $K$, magnitudes, end-point, seeds, released artefacts) & -- & release probes, directions and control seeds \\
\bottomrule
\end{tabular}
\end{table*}
